
\documentclass[10pt]{article}
\usepackage[T1]{fontenc}
\usepackage[utf8]{inputenc}
\usepackage{lmodern}
\usepackage{amsmath,amssymb,mathtools}
\usepackage{geometry}
\usepackage{enumitem}
\usepackage{tabularx}
\usepackage{tikz-cd}
\geometry{margin=0.86in}
\setlength{\parindent}{1.2em}
\setlength{\parskip}{0.12em}
\emergencystretch=3em
\setlist{leftmargin=2.1em,itemsep=0.08em,topsep=0.2em}
\newcommand{\C}{\mathbb C}
\newcommand{\Q}{\mathbb Q}
\newcommand{\Z}{\mathbb Z}
\newcommand{\F}{\mathbb F}
\newcommand{\Ql}{\mathbb Q_\ell}
\newcommand{\Zl}{\mathbb Z_\ell}
\newcommand{\A}{\mathbb A}
\newcommand{\G}{\mathbb G}
\newcommand{\Gal}{\operatorname{Gal}}
\newcommand{\GL}{\operatorname{GL}}
\newcommand{\Hom}{\operatorname{Hom}}
\newcommand{\Ind}{\operatorname{Ind}}
\newcommand{\Res}{\operatorname{Res}}
\newcommand{\Inf}{\operatorname{Inf}}
\newcommand{\detm}{\operatorname{det}}
\newcommand{\dimv}{\operatorname{dim}}
\newcommand{\Tr}{\operatorname{Tr}}
\newcommand{\Fr}{\operatorname{Frob}}
\newcommand{\rk}{\operatorname{rk}}
\newcommand{\ppcm}{\operatorname{ppcm}}
\newcommand{\Spec}{\operatorname{Spec}}
\newcommand{\iso}{\xrightarrow{\sim}}
\newcommand{\into}{\hookrightarrow}
\newcommand{\onto}{\twoheadrightarrow}

\providecommand{\Ker}{\operatorname{Ker}}
\providecommand{\coker}{\operatorname{coker}}
\providecommand{\Stab}{\operatorname{Stab}}
\providecommand{\Char}{\operatorname{Hom}}
\providecommand{\Frac}{\operatorname{Frac}}
\providecommand{\ord}{\operatorname{ord}}
\providecommand{\im}{\operatorname{im}}


\providecommand{\OO}{\mathcal O}
\providecommand{\mm}{\mathfrak m}
\providecommand{\ab}{\mathrm{ab}}
\providecommand{\geom}{\mathrm{geom}}
\providecommand{\cD}{\mathcal D}
\providecommand{\cF}{\mathcal F}
\providecommand{\Frob}{\operatorname{Frob}}
\providecommand{\vol}{\operatorname{vol}}


\providecommand{\R}{\mathbb R}
\providecommand{\OO}{\mathcal O}
\providecommand{\mm}{\mathfrak m}
\providecommand{\ab}{\mathrm{ab}}
\providecommand{\geom}{\mathrm{geom}}
\providecommand{\cD}{\mathcal D}
\providecommand{\cF}{\mathcal F}
\providecommand{\Frob}{\operatorname{Frob}}
\providecommand{\vol}{\operatorname{vol}}

\begin{document}

\begin{center}
{\Large\bfseries LES CONSTANTES DES ÉQUATIONS FONCTIONNELLES DES FONCTIONS \(L\)}\\[0.4em]
{\normalsize 2.2--3.8}
\end{center}


Pour \(s\in\C\), on note \(\omega_s\) le quasi-caractère
\[
x\longmapsto \|x\|^s
\]
de \(K^*\) dans \(\C^*\).

Si \(K\) est non archimédien, c'est-à-dire \(K\ne\R,\C\), on note \(\mathcal O\) l'anneau de la valuation \(v\) de \(K\), \(\pi\) une uniformisante, \(k\) le corps résiduel \(\mathcal O/(\pi)\), \(p\) la caractéristique de \(k\), et on pose
\[
d=[k:\F_p],\qquad q=p^d=\#k.
\]
On a donc
\[
\|x\|=q^{-v(x)}.
\]

\paragraph{2.2.2.} Supposons \(K\) non archimédien. On désigne par \(\overline K\) une clôture séparable de \(K\), et on note \(\overline{\mathcal O}\) la clôture intégrale de \(\mathcal O\) dans \(\overline K\), \(\overline k\) le corps résiduel de \(\overline{\mathcal O}\), donc une clôture algébrique de \(k\), et \(v\) la valuation de \(\overline K\), à valeurs dans \(\Q\), qui prolonge celle de \(K\).

On dispose d'une filtration en trois crans du groupe de Galois
\[
\Gal(\overline K/K)\supset I\supset P
\]
de quotients successifs
\[
\begin{array}{rcl}
\Gal(\overline K/K)/I&\iso&\Gal(\overline k/k)\iso\widehat\Z,\\[0.35em]
I/P&\iso&\widehat\Z(1)(\overline k),\\[0.35em]
P&:&\text{un pro-\(p\)-groupe.}
\end{array}
\]
On l'obtient comme suit; pour les démonstrations, on renvoie à [10].

\begin{enumerate}[label=\alph*)]
\item Par transport de structure, \(\Gal(\overline K/K)\) agit sur \(\overline{\mathcal O}\) et \(\overline k\). Le groupe d'inertie \(I\) est le noyau de la flèche de réduction
\[
v':\Gal(\overline K/K)\longrightarrow \Gal(\overline k/k).
\]
Le quotient \(\Gal(\overline k/k)\) est canoniquement isomorphe à \(\widehat\Z\), de générateur topologique la substitution de Frobenius
\[
\varphi:x\longmapsto x^q.
\]

\item Le groupe d'inertie sauvage \(P\) est l'unique pro-\(p\)-groupe de Sylow de \(I\). Le morphisme équivariant
\[
t:I\longrightarrow \widehat\Z(1)(\overline k),
\]
de noyau \(P\), est caractérisé par la congruence, modulo l'idéal maximal de \(\overline{\mathcal O}\),
\[
\frac{\sigma(x)}{x}=t(\sigma)\,v(x)
\qquad(\sigma\in I,\ x\in\overline K).
\]
Pour \(\ell\ne p\) premier, on note \(t_\ell\) l'application composée
\[
I\longrightarrow\widehat\Z(1)(\overline k)\longrightarrow\Z_\ell(1)(\overline k).
\]
\end{enumerate}

\paragraph{2.2.3.} Soit \(k\) un corps fini à \(q\) éléments, de clôture algébrique \(\overline k\). On note \(W(\overline k/k)\) le sous-groupe de \(\Gal(\overline k/k)\simeq\widehat\Z\) engendré par
\[
\varphi:x\longmapsto x^q.
\]
On a canoniquement
\[
W(\overline k/k)\simeq\Z
\]
engendré par \(\varphi\). Le Frobenius géométrique \(F\in W(\overline k/k)\) est par définition \(\varphi^{-1}\).

\paragraph{2.2.4.} Soient \(K,\overline K\) comme en 2.2.2. Le groupe de Weil \(W(\overline K/K)\), ou groupe de Weil absolu de \(K\), est le sous-groupe de \(\Gal(\overline K/K)\) formé des \(\sigma\) tels que \(v'(\sigma)\) soit une puissance entière du Frobenius \(\varphi\). On le munit de la topologie induisant la topologie naturelle de \(I\), et pour laquelle \(I\) est ouvert. Le groupe \(\Gal(\overline K/K)\) est le complété de \(W(\overline K/K)\) pour la topologie des sous-groupes ouverts d'indice fini; pour \(L\) une extension finie de \(K\) dans \(\overline K\), le sous-groupe correspondant de \(W(\overline K/K)\) s'identifie à \(W(\overline K/L)\). On appellera parfois Frobenius géométriques les éléments de \(\Gal(\overline K/K)\) ou \(W(\overline K/K)\) d'image \(F\).

\paragraph{2.2.5.} Supposons \(K\) archimédien, de clôture algébrique \(\overline K\). On définit \(W(\overline K/K)\) comme l'extension suivante de \(\Gal(\overline K/K)\) par \(\overline K^*\):
\begin{enumerate}[label=\arabic*)]
\item \(K\simeq\R\): \(W(\overline K/K)\) est engendré par \(\overline K^*\) et un élément \(F\), avec
\[
F^2=-1,\qquad FzF^{-1}=\overline z\quad(z\in\overline K^*).
\]
\item \(K\simeq\C\): \(W(\overline K/K)=\overline K^*\).
\end{enumerate}

\paragraph{2.3.} La théorie du corps de classe local fournit un isomorphisme entre \(W(\overline K/K)^{\ab}\) et \(K^*\).

Supposons \(K\) non archimédien. Pour une raison expliquée en 3.6, nous normaliserons alors cet isomorphisme, c'est-à-dire choisirons lui ou son inverse, de telle sorte que les Frobenius géométriques correspondent aux uniformisantes:
\[
\begin{tikzcd}[column sep=2.7em,row sep=1.9em]
P \arrow[r,hook] \arrow[d] & I \arrow[r,hook] \arrow[d] &
W(\overline K/K) \arrow[r] \arrow[d] &
W(\overline k/k)\simeq\Z \arrow[d,"{-1}"]\\
1+(\pi) \arrow[r,hook] & \mathcal O^* \arrow[r,hook] &
K^* \arrow[r,"v"] & \Z .
\end{tikzcd}
\]
Une représentation de \(W(\overline K/K)\) est non ramifiée, respectivement modérément ramifiée, si elle est triviale sur \(I\), respectivement sur \(P\). De même, un quasi-caractère \(\chi\) de \(K^*\) est non ramifié, respectivement modérément ramifié, s'il est trivial sur \(\mathcal O^*\), respectivement sur \(1+(\pi)\). Pour \(\ell\) premier à \(p\), l'action de \(\Gal(\overline K/K)\) sur \(\Z_\ell(1)\) est non ramifiée, décrite par le quasi-caractère
\[
\omega_1=\|x\|
\]
de \(K^*\).

Pour \(K\) complexe, l'isomorphisme de la théorie du corps de classe local est celui évident sur 2.2.5; pour \(K\) réel, c'est celui qui rend vrais 2.3.1 et 2.3.2 ci-dessous.

Soit \(L\subset\overline K\) une extension finie de \(K\), de sorte que \(W(\overline K/L)\subset W(\overline K/K)\). Les diagrammes suivants, où \(N\) est la norme et \(t\) le transfert, sont commutatifs ([10], XIII, §4):
\[
\tag{2.3.1}
\begin{tikzcd}[column sep=3.3em,row sep=2.2em]
W(\overline K/L) \arrow[r] \arrow[d] &
W(\overline K/L)^{\ab} \arrow[r,"\sim"] \arrow[d] &
L^* \arrow[d,"N_{L/K}"]\\
W(\overline K/K) \arrow[r] &
W(\overline K/K)^{\ab} \arrow[r,"\sim"] &
K^* ,
\end{tikzcd}
\]
\[
\tag{2.3.2}
\begin{tikzcd}[column sep=3.5em,row sep=2.3em]
W(\overline K/K)^{\ab} \arrow[r,"\sim"] \arrow[d,"t"'] &
K^* \arrow[d]\\
W(\overline K/L)^{\ab} \arrow[r,"\sim"] &
L^* .
\end{tikzcd}
\]

\paragraph{2.4.} Soient \(K\) un corps de fonctions d'une variable sur \(\F_p\) et \(k\) son corps des constantes, c'est-à-dire la fermeture algébrique de \(\F_p\) dans \(K\). Soient \(\overline K\) une clôture séparable de \(K\), \(\overline k\) la clôture algébrique de \(k\) dans \(\overline K\), et \(\Gal^0(\overline K/K)\) le noyau de l'application de restriction de \(\Gal(\overline K/K)\) dans \(\Gal(\overline k/k)\). Le groupe de Weil global absolu \(W(\overline K/K)\) est, en analogie avec 2.2.4, défini par le diagramme
\[
\begin{tikzcd}[column sep=1.7em,row sep=2.1em]
0 \arrow[r] & \Gal^0(\overline K/K) \arrow[r] \arrow[d,equal] &
W(\overline K/K) \arrow[r] \arrow[d,hook] &
W(\overline k/k) \arrow[r] \arrow[d,hook] & 0\\
0 \arrow[r] & \Gal^0(\overline K/K) \arrow[r] &
\Gal(\overline K/K) \arrow[r] &
\Gal(\overline k/k) \arrow[r] & 0 .
\end{tikzcd}
\]
Soient \(v\) une place de \(K\), et supposons choisie une place de \(\overline K\) au-dessus de \(v\). On dispose alors d'un diagramme commutatif
\[
\begin{tikzcd}[column sep=1.6em,row sep=2.1em]
0 \arrow[r] & I_v \arrow[r] \arrow[d,hook] &
W(\overline K_v/K_v) \arrow[r] \arrow[d,hook] &
W(\overline k_v/k_v) \arrow[r] \arrow[d,hook] & 0\\
0 \arrow[r] & \Gal^0(\overline K/K) \arrow[r] &
W(\overline K/K) \arrow[r] &
W(\overline k/k) \arrow[r] & 0 .
\end{tikzcd}
\]
La théorie du corps de classe global fournit un isomorphisme
\[
W(\overline K/K)^{\ab}\iso\A_K^*/K^*
\]
où \(\A_K^*/K^*\) est le groupe des classes d'idèles, rendant commutatifs les diagrammes
\[
\begin{tikzcd}[column sep=3.0em,row sep=2.0em]
W(\overline K_v/K_v)^{\ab} \arrow[r] \arrow[d] &
W(\overline K/K)^{\ab} \arrow[d]\\
K_v^* \arrow[r] & \A_K^*/K^* .
\end{tikzcd}
\]

\paragraph{2.5.} Pour la définition des groupes de Weil globaux dans le cas des corps de nombres, on renvoie à Weil [19] ou à [2], XIV. Nous ne nous en servirons qu'incidemment.

\section*{3. Rappels sur les fonctions \(L\)}

\paragraph{3.1.} Soit \(K\) un corps local et reprenons les notations 2.2.1. On désignera par \(dx\) une mesure de Haar sur \(K\), par \(d^*x\) une mesure de Haar sur \(K^*\), et par \(\psi\) un caractère additif non trivial de \(K\).

\paragraph{3.2.} Pour un quasi-caractère, c'est-à-dire un homomorphisme continu,
\[
\chi:K^*\longrightarrow\C^*,
\]
on définit comme suit \(L(\chi)\in\C\cup\{\infty\}\).

\[
\tag{3.2.1}
K\simeq\R.
\]
On pose
\[
\Gamma_\R(s)=\pi^{-s/2}\Gamma(s/2),
\]
et, pour \(x\) le plongement de \(K\) dans \(\C\) et \(N=0\) ou \(1\),
\[
L(x^{-N}\omega_s)=\Gamma_\R(s).
\]

\[
\tag{3.2.2}
K\simeq\C.
\]
On pose
\[
\Gamma_\C(s)=2(2\pi)^{-s}\Gamma(s),
\]
et, pour \(z\) un plongement de \(K\) dans \(\C\) et \(N\ge0\),
\[
L(z^{-N}\omega_s)=\Gamma_\C(s).
\]

\[
\tag{3.2.3}
K\text{ non archimédien.}
\]
On pose
\[
L(\chi)=1\quad(\chi\text{ ramifié}),
\qquad
L(\chi)=\frac1{1-\chi(\pi)}\quad(\chi\text{ non ramifié}).
\]

\paragraph{3.3.} \(\psi\) et \(dx\) étant donnés, on pose, pour \(f\) une fonction \(C^\infty\) à support compact sur \(K\),
\[
\widehat f(y)=\int_K f(x)\psi(xy)\,dx.
\]
On définit \(\varepsilon(\chi,\psi,dx)\in\C^*\) par l'équation fonctionnelle locale de Tate, indépendante de \(f\),
\[
\tag{3.3.1}
\frac{\displaystyle\int_{K^*}\widehat f(x)\,\omega_1\chi^{-1}(x)\,d^*x}
     {L(\omega_1\chi^{-1})}
=
\varepsilon(\chi,\psi,dx)\,
\frac{\displaystyle\int_{K^*}f(x)\chi(x)\,d^*x}
     {L(\chi)}.
\]
On prend la même mesure de Haar \(d^*x\) dans les deux membres. Les deux membres sont définis par prolongement analytique en \(\chi\) dans les familles \(\chi\omega_s\), \(s\in\C\). Dans ces familles, \(\varepsilon\) est de la forme \(Ae^{Bs}\).

On a
\[
\tag{3.3.2}
\varepsilon(\chi,\psi,a\,dx)=a\,\varepsilon(\chi,\psi,dx),
\]
\[
\tag{3.3.3}
\varepsilon(\chi,\psi(ax),dx)=\chi(a)\|a\|^{-1}\varepsilon(\chi,\psi,dx).
\]

\paragraph{3.4.} Pour \(K\) non archimédien, nous poserons:
\[
n(\psi)=\text{ le plus grand entier }n\text{ tel que }\psi|_{\pi^{-n}\mathcal O}=1;
\]
\[
a(\chi)=
\begin{cases}
0,&\chi\text{ non ramifié},\\
\text{le plus petit }m\text{ tel que }\chi|_{1+(\pi)^m}=1,&\chi\text{ ramifié};
\end{cases}
\]
\[
sw(\chi)=
\begin{cases}
0,&\chi\text{ non ou modérément ramifié},\\
a(\chi)-1,&\chi\text{ ramifié};
\end{cases}
\]
et \(\gamma\) désignera un élément de \(K^*\) de valuation
\[
v(\gamma)=n(\psi)+sw(\chi)+1.
\]
La constante locale \(\varepsilon\) est donnée par (3.3.2), (3.3.3), et les formules suivantes.

\[
\tag{3.4.1}
K\simeq\R.
\]
Soient \(x\) le plongement de \(K\) dans \(\C\) et \(N=0\) ou \(1\). Pour \(\psi=\exp(2\pi ix)\) et \(dx\) la mesure de Lebesgue,
\[
\varepsilon(x^{-N}\omega_s,\psi,dx)=i^N.
\]

\[
\tag{3.4.2}
K\simeq\C.
\]
Soient \(z\) un plongement de \(K\) dans \(\C\) et \(N\ge0\). Pour
\[
\psi=\exp(2\pi i\,\Tr_{\C/\R}(z))
\]
et
\[
dx=|dz\wedge d\overline z|=2\,da\,db\quad(z=a+bi),
\]
on a
\[
\varepsilon(z^{-N}\omega_s,\psi,dx)=i^N.
\]

\paragraph{3.4.3.} \(K\) non archimédien. Si \(\chi\) est non ramifié,
\[
\int_{\mathcal O}dx=1,\qquad n(\psi)=0,
\]
alors
\[
\tag{3.4.3.1}
\varepsilon(\chi,\psi,dx)=1.
\]
Pour \(\chi\) ramifié, on a
\[
\tag{3.4.3.2}
\varepsilon(\chi,\psi,dx)=
\int_{K^*}\chi^{-1}(x)\psi(x)\,dx
=
\sum_n\int_{v(x)=n}\chi^{-1}(x)\psi(x)\,dx
=
\int_{\gamma^{-1}\mathcal O^*}\chi^{-1}(x)\psi(x)\,dx.
\]

On verra qu'il est naturel d'unifier ces deux formules de la façon suivante. Si on définit \(\varepsilon_0(\chi,\psi,dx)\) par
\[
\varepsilon(\chi,\psi,dx)=\varepsilon_0(\chi,\psi,dx)
\quad(\chi\text{ ramifié}),
\]
et
\[
\tag{3.4.3.3}
\varepsilon(\chi,\psi,dx)=\varepsilon_0(\chi,\psi,dx)\,\chi(\pi)^{-1}
\quad(\chi\text{ non ramifié}),
\]
on a
\[
\tag{3.4.3.4}
\varepsilon_0(\chi,\psi,dx)=
\int_{\gamma^{-1}\mathcal O^*}\chi^{-1}(x)\psi(x)\,dx.
\]
De (3.4.3.3) et (3.4.3.4), on déduit que, pour \(\omega\) non ramifié,
\[
\tag{3.4.3.5}
\varepsilon_0(\chi\omega,\psi,dx)
=
\omega(\pi)^{n(\psi)+sw(\chi)+1}\varepsilon_0(\chi,\psi,dx),
\]
et
\[
\varepsilon(\chi\omega,\psi,dx)
=
\omega(\pi)^{n(\psi)+a(\chi)}\varepsilon(\chi,\psi,dx).
\]

\paragraph{3.5.} Supposons \(K\) non archimédien, et reprenons les notations 2.2.1. Soit \(V\) un espace vectoriel de dimension finie sur un corps \(E\) de caractéristique \(0\), ou sur \(\C\). Les représentations
\[
\rho:W(\overline K/K)\longrightarrow\GL(V)
\]
seront toujours supposées continues: un sous-groupe ouvert de \(I\) agit trivialement. Pour \(V\) l'espace d'une représentation et \(V^I\) le sous-espace des invariants sous l'inertie, on pose
\[
\tag{3.5.1}
Z(V;t)=\det(1-Ft^d,V^I)^{-1},
\]
\[
\tag{3.5.2}
L(V)=Z(V;1)\qquad(E\in E^*\cup\{\infty\}).
\]
Notons
\[
\omega^t:W(\overline K/K)\longrightarrow E((t))^*
\]
le caractère non ramifié tel que \(\omega^t(F)=t^d\). En un sens évident, on a \(\omega_s=\omega^{p^{-s}}\). On a
\[
\tag{3.5.3}
Z(V,t)=L(V\otimes\omega^t).
\]

\paragraph{3.6.} C'est pour que (3.2.3) et (3.5.2) soient compatibles que nous avons normalisé l'isomorphisme du corps de classe local de telle sorte que les Frobenius géométriques correspondent aux uniformisantes. Tant (3.2.3) que (3.5.2) sont très naturels: 3.2.3 s'impose du point de vue analytique (cf. 3.3.1), et 3.5.2 suit la tradition des géomètres; c'est le Frobenius géométrique qui tend à agir sur la cohomologie des variétés algébriques avec pour valeurs propres des entiers algébriques.

\paragraph{3.7.} Les représentations continues complexes irréductibles de \(W(\overline K/K)\) sont, pour \(K\) complexe, les quasi-caractères. Pour \(K\) réel, elles sont, outre les quasi-caractères de \(K^*\), les représentations induites par les quasi-caractères \(\chi\) de \(\overline K^*\) tels que \(\chi\ne\chi\circ\sigma\), pour \(\sigma\) la conjugaison complexe de \(\overline K\).

Pour toute représentation complexe \(V\), exprimée, dans le groupe de Grothendieck des représentations de \(W(\overline K/K)\), comme somme de représentations simples, on définit \(L(V)\) comme suit:
\[
\tag{3.7.1}
K\simeq\R,\quad
V=\sum_i[\chi_i]+\sum_j\Ind_{\overline K^*}^{W}(\chi_j),
\]
\[
L(V)=\prod_i L_K(\chi_i)\prod_j L_{\overline K}(\chi_j).
\]
\[
\tag{3.7.2}
K\simeq\C,\quad V=\sum_i[\chi_i],
\qquad
L(V)=\prod_iL(\chi_i).
\]

\paragraph{Proposition 3.8.} Soient \(K\) un corps local et \(\overline K\) une clôture algébrique de \(K\).
\begin{enumerate}[label=(\roman*)]
\item Pour toute suite exacte
\[
0\longrightarrow V'\longrightarrow V\longrightarrow V''\longrightarrow0
\]
de représentations complexes de \(W(\overline K/K)\), on a
\[
\tag{3.8.1}
L(V)=L(V')L(V'').
\]
\item Soient \(L\) une extension finie séparable de \(K\) dans \(\overline K\), et \(V_L\) une représentation complexe de \(W(\overline K/L)\). Soit \(V_K\) la représentation induite de \(W(\overline K/K)\). On a
\[
\tag{3.8.2}
L(V_K)=L(V_L).
\]
\end{enumerate}

L'assertion (i) est triviale. Dans le cas archimédien, l'assertion (ii) résulte de la formule de duplication
\[
\Gamma_\C(s)=\Gamma_\R(s)\Gamma_\R(s+1),
\]
correspondant à
\[
\Ind([\omega_s])=[\omega_s]+[\chi^{-1}\omega_{s+1}].
\]

Prouvons (ii) pour \(K\) non archimédien. Notant \(k\) et \(\ell\) les corps résiduels de \(K\) et \(L\), on a
\[
V_K^{I_K}=
\left(\Ind_{W(\overline K/L)}^{W(\overline K/K)}V_L\right)^{I_K}
\simeq
\Ind_{W(\overline\ell/\ell)}^{W(\overline k/k)}(V_L^{I_L})
\]
comme \(W(\overline k/k)\)-représentations, les deux membres s'identifiant à
\[
\Hom_{W(\overline K/L)}(W(\overline k/k),V_L),
\]
fonctions covariantes de \(W(\overline k/k)\) dans \(V_L\), \(W(\overline K/L)\) agissant sur \(W(\overline k/k)\) par translation. Cf. le diagramme commutatif

\end{document}
